% TOP_PROBLEM: {"id": 4700032, "title": "The 196 conjecture", "category_id": 2, "category": {"id": 2, "name": "combinatorics", "display_name": "Combinatorics", "description": "Counting problems, graph theory, discrete structures.", "slug": "combinatorics", "order_index": 2, "created_at": "2026-07-31T15:26:25.671Z"}, "status": "open", "problem_number": "AMR-046-0032", "set_id": 13}
% FIRST_POSTED: 2026-10-02
% ENTRY_KIND: statement_only
% UnsolvedMath ID 4700032; source catalogue code AMR-046-0032
% !TeX program = lualatex
% Definition and sources only; this file is not an LLM solution attempt.
\documentclass[11pt,a4paper]{article}
\usepackage[margin=25mm]{geometry}
\usepackage{fontspec}
\setmainfont{lmroman10-regular.otf}[BoldFont=lmroman10-bold.otf,
 ItalicFont=lmroman10-italic.otf,BoldItalicFont=lmroman10-bolditalic.otf]
\usepackage{amsmath,amssymb,mathtools}
\usepackage{unicode-math}
\setmathfont{latinmodern-math.otf}
\AtBeginDocument{\let\setminus\smallsetminus}
\usepackage{xurl,hyperref}
\hypersetup{colorlinks=true,urlcolor=blue}
\setcounter{secnumdepth}{0}
\setlength{\parindent}{0pt}
\setlength{\parskip}{5pt}
\setlength{\emergencystretch}{3em}
\begin{document}
\section*{The 196 conjecture}\label{4700032}
{\small UnsolvedMath ID 4700032 · AMR-046-0032 · Combinatorics · AMR Open Problem Lists}
\subsection{Definitions and mathematical statement}
For a positive integer $n$ with decimal digits $d_{r-1}\cdots d_0$, define
$\operatorname{rev}(n)$ to be the integer represented by the reversed
string $d_0\cdots d_{r-1}$; any resulting leading zeros are discarded.
Set $R(n)=n+\operatorname{rev}(n)$ and form the sequence
$n_0=n$, $n_{j+1}=R(n_j)$. A positive integer is a palindrome when its
decimal digit string reads the same forwards and backwards.

Define
\[
 \mathcal L=\{n\ge1:n_j\text{ is not a palindrome for every }j\ge1\}.
\]
The source asks two questions: is $\mathcal L$ infinite, and is its least
element equal to $196$? The second statement includes both
$196\in\mathcal L$ and the assertion that every positive integer smaller
than $196$ reaches a palindrome after finitely many positive iterations.

The requirement is avoidance forever. A large number of computed
nonpalindromic iterates does not certify membership in $\mathcal L$.
The positive-iteration convention is the one in the source: the initial
term $n_0$ is not itself tested. Decimal notation is fixed throughout,
and neither reversal in another base nor retaining artificial leading
zeros changes the operation under consideration. Different starting
integers can eventually enter the same orbit, but the infinitude question
counts starting integers, not distinct tail-equivalence classes of
orbits. Any proof of the claimed minimum must in particular establish
nontermination of the reverse-and-add process starting at $196$.
\subsection{Short English statement}
Are infinitely many decimal reverse-and-add sequences forever nonpalindromic, and is $196$ the smallest starting value?
\subsection{Sources}
\begin{itemize}
\item[S1] ULAM.AI and the UnsolvedMath Contributors, supplied problems.json, record 4700032 (AMR-046-0032), AMR Open Problem Lists. Catalogue identity and source transcription; CC BY 4.0.
\url{https://huggingface.co/datasets/ulamai/UnsolvedMath}.
\item[S2] Gasull - Some open problems in low dimensional dynamical systems (2020), Problem environment 32: The 196 conjecture. Original source indexed by AMR-046-0032.
\url{https://arxiv.org/abs/2012.02524}.
\end{itemize}
\end{document}
